\section*{摘要}

将无人机航拍图像上传至边缘视觉--语言模型（VLM）可实现自动视觉问答，
但对每个问题均传输图像并执行模型推理会产生较高的通信与计算能耗。
针对这一问题，本文提出一种面向结构化航拍视觉问答的证据级语义通信系统，
根据问题的信息需求联合选择传输内容与接收端处理方式。
发送端在检测器生成的结构化目标信息与图像之间进行选择，
接收端分别采用轻量符号解码器与冻结的 VLM 生成答案。
利用不同问题类型对两类证据的需求差异，设计无需训练的类型选择规则，
使可由目标信息回答的问题避免图像传输与 VLM 推理。
在此基础上，构建结合逐样本检测信息的学习型路由器，
并引入能耗价格调节准确率与能耗之间的权衡。
实验使用两个航拍数据集，并区分传输、机载检测与边缘推理的能耗。
在 VisDrone 的三种信道模型下，类型规则相较于固定图像传输提高了准确率，
并将核算的系统能耗降低约 $55\%$。
在 DroneVehicle 上，通过验证集选择能耗权重的学习型路由器，
以每题 $4.47$ J 的能耗达到 $74.60\%$ 的准确率；
类型规则的对应结果为 $11.21$ J 和 $72.93\%$。
匹配接收器的比较表明，针对接收器进行适配可以带来收益，
但非线性路由并非始终优于线性基线。
这些结果支持根据任务联合选择证据与回答计算方式。
节能收益针对所定义的无人机与边缘端合计能耗模型，不代表实测无人机续航收益。

\par\medskip
\noindent\textbf{关键词：}能量效率、证据级路由、语义通信、无人机、视觉--语言模型、视觉问答。
